\section{5.1 Runtime Environment}
\begin{itemize}
  \item OS: Ubuntu 22.04
  \item Python: 3.10 (\texttt{venv})
  \item Main libraries: OpenCV, NumPy, Ultralytics, MMPose/MMEngine
  \item Camera hardware: Hikvision DS-E12 USB cameras
  \item Runtime resolution and frame rate: 1280x720 at 15 FPS target
\end{itemize}

\section{5.2 Configuration Layer}
Two camera configs co-exist in repository:
\begin{itemize}
  \item \texttt{config/cameras.yaml} (older generic mapping),
  \item \texttt{garage\_lab\_combined/config/cameras.yaml} (role-based mapping with by-path links).
\end{itemize}

Role-based mapping (\texttt{camNorth}, \texttt{camEast}, \texttt{camSouth}, \texttt{camWest}) is necessary because control and extrinsics logic depend on physical orientation, not on changing \texttt{/dev/video*} indices.

\section{5.3 Intrinsics Capture Automation}
\texttt{auto\_capture\_charuco\_multi.py} was improved to auto-save frames when sufficient corners are detected, replacing fixed delay logic. This addressed practical constraints where the operator was physically far from the workstation.

A preview window and corner-threshold trigger improved data quality and reduced unusable captures.

\section{5.4 Intrinsics Solving}
\texttt{calibrate\_intrinsics\_from\_images.py} handles per-camera image sets and outputs JSON files with:
\begin{itemize}
  \item camera matrix,
  \item distortion coefficients,
  \item reprojection error,
  \item frames used.
\end{itemize}

Practical lesson: an A4 full-page ChArUco board with accurate printed dimensions and high contrast improved corner distribution quality over earlier board setups.

\section{5.5 Extrinsics Solving and Iterative Hardening}
\texttt{calibrate\_extrinsics\_apriltag\_robust.py} supports robust fitting and camera-specific tag maps. This was essential when cameras were moved and visibility changed.

Engineering decisions that improved results:
\begin{itemize}
  \item re-measuring camera coordinates in world frame,
  \item correcting swapped east/west camera positions in \texttt{Dimensions.txt},
  \item updating moved AprilTag coordinates (e.g., IDs 13, 14, 18, 21, 22),
  \item using two-stage optimization (strict tags first, expanded tags second).
\end{itemize}

\section{5.6 Overlay Validation}
\texttt{validate\_extrinsics\_overlay.py} provides red/green projection diagnostics:
\begin{itemize}
  \item green: detected tag corners,
  \item red: expected projected corners from calibration model.
\end{itemize}

Overlay checks prevented silent geometric failures and guided targeted correction of tags/camera poses.

\section{5.7 Capture and Session Management}
\texttt{record\_short\_clips\_multi.py} supports short synchronized clips with keyboard controls. For joint tests, \texttt{auto\_record\_joint\_trials.py} was adapted to include operator settle time after pressing \texttt{r}, making single-operator experiments feasible.

Session folder conventions (\texttt{clips}, \texttt{results}, \texttt{reports}, \texttt{visualizations}) simplified reproducibility and report generation.

\section{5.8 3D Processing Pipeline Decisions}
\texttt{process\_4cam\_to\_3d.py} options used in final experiments:
\begin{itemize}
  \item \texttt{--ball-min-cams 2} (practical minimum under visibility constraints),
  \item reprojection threshold gating,
  \item optional EMA and speed gates,
  \item \texttt{--pose-min-cams 3} for stricter joint reliability in main smoke runs.
\end{itemize}

This parameterization reflects a tradeoff: stricter rules improve quality but reduce frame coverage.

\section{5.9 Rendering and Presentation}
\texttt{render\_arena\_ball\_skeleton.py} was tuned for stakeholder-friendly visuals:
\begin{itemize}
  \item arena wireframe and camera arrows,
  \item black skeleton lines for contrast,
  \item ball trajectory with visible markers,
  \item optional no-smoothing render for diagnostic honesty.
\end{itemize}

Additional scripts (\texttt{render\_multiviews.py}, live view) supported top/side perspective analysis.

\section{5.10 CPU/GPU Split Discussion}
The current pipeline can run CPU-only but with limited real-time throughput. A practical split strategy is:
\begin{itemize}
  \item GPU: ball detector + pose model inference,
  \item CPU: decode, synchronization logic, triangulation, rendering control.
\end{itemize}

This separation reduces contention and is suitable for future live deployment goals.

\section{5.11 Failure Modes and Mitigations}
Observed failure modes:
\begin{itemize}
  \item camera index remapping after reconnects,
  \item one-camera movement causing global 3D drift,
  \item multiple-person identity switching,
  \item ball outliers during fast motion,
  \item low overlap causing 2-camera-only triangulation.
\end{itemize}

Mitigations applied:
\begin{itemize}
  \item by-path camera mapping,
  \item partial camera recalibration and merge,
  \item single-subject protocol enforcement,
  \item reprojection/speed gates,
  \item redesigned GT point placement for >=3-camera visibility where possible.
\end{itemize}

\section{5.12 Why This Implementation Is Research-Ready}
The implementation is not only script-complete; it is protocolized. Every critical phase has:
\begin{itemize}
  \item data artifacts,
  \item configuration snapshots,
  \item metrics outputs,
  \item visual diagnostics.
\end{itemize}

This traceability enables scientific defense and iterative improvement.


\section{5.13 Multi-Session Calibration Management}
A major practical difficulty was the accumulation of many calibration files across dates and variants. To prevent accidental use of outdated calibration, the workflow now follows:
\begin{enumerate}
  \item create dated calibration IDs,
  \item run overlay validation snapshots,
  \item pin an explicit \texttt{extrinsics\_final\_<timestamp>.json} for each evaluation session,
  \item record this file path in experiment logs.
\end{enumerate}

This discipline reduced silent failure cases where processing succeeded but used inconsistent geometry.

\section{5.14 Camera Movement Incident Handling}
The project experienced multiple real incidents where one or more cameras moved. The adopted response procedure:
\begin{itemize}
  \item identify moved camera(s),
  \item capture short recalibration clip,
  \item extract stills,
  \item rerun robust extrinsics for affected camera(s),
  \item merge updated camera entries with stable entries,
  \item validate with overlays,
  \item rerun smoke processing before any GT campaign.
\end{itemize}

This procedure is now part of operational knowledge and should be retained in final project documentation.

\section{5.15 Runtime Performance and Throughput Strategy}
At 1280x720 and 4 cameras, the practical throughput depends on inference backend and scene complexity. For future live usage:
\begin{itemize}
  \item decouple capture and inference queues,
  \item keep deterministic frame indexing,
  \item prioritize inference on GPU,
  \item run triangulation and logging on CPU,
  \item use adaptive frame skipping only when confidence remains stable.
\end{itemize}

The objective is not maximum FPS but stable, interpretable, and safe timing behavior.

\section{5.16 Visualization as a Debug Instrument}
Several major issues were discovered through visualization first:
\begin{itemize}
  \item skeleton appearing outside plausible room area,
  \item ball vanishing due to over-strict filtering,
  \item coordinate squeezes when wrong extrinsics were used,
  \item person entry confusion when another person stayed in scene.
\end{itemize}

Therefore, visual outputs are integrated into validation pipeline, not only into presentation pipeline.

\section{5.17 Engineering Trade-Off Summary}
Key trade-offs made in implementation:
\begin{itemize}
  \item strict thresholds improve reliability but reduce detection coverage,
  \item smoothing reduces jitter but can hide dynamics,
  \item partial recalibration is faster but may preserve old global inconsistencies,
  \item high resolution improves detail but reduces real-time throughput.
\end{itemize}

These trade-offs are documented explicitly so future development can adjust based on application priority (accuracy, latency, safety, or ease of operation).



\section{5.18 Detailed Command-Line Pipeline for Reproducibility}
A complete session generally follows this order:
\begin{enumerate}
  \item record synchronized clips,
  \item run calibration updates if needed,
  \item validate overlays,
  \item process 4-camera clips to 3D JSON,
  \item render arena outputs,
  \item run GT evaluation scripts,
  \item review summary metrics and diagnostics.
\end{enumerate}

Each stage writes durable outputs. No stage depends on hidden in-memory state, which improves reproducibility.

\section{5.19 File Versioning and Non-Overwrite Policy}
The workflow intentionally avoids silent overwrite of major artifacts by:
\begin{itemize}
  \item timestamped output paths,
  \item suffixes like \texttt{\_raw}, \texttt{\_fixed}, \texttt{\_corrected},
  \item per-session roots with subfolders.
\end{itemize}

This design made comparative analysis feasible (e.g., old vs new extrinsics) and protected against accidental data loss.

\section{5.20 Engineering of Manual Flash Synchronization Workflow}
Manual synchronization with flashlight was retained because it is transparent and operator-verifiable. Automated offset estimates were available but less trusted under varying lighting and frame truncation artifacts. The final process used manual frame identification and explicit clip trimming for deterministic alignment.

\section{5.21 Integration with Legacy Renderers}
Legacy renderers in \texttt{src/core} were preserved to cross-check the new renderer's behavior. If one renderer showed implausible behavior while another did not, this indicated differences in coordinate assumptions or smoothing logic. Cross-render consistency checks improved confidence in final outputs.

\section{5.22 Practical Hardware Notes}
Observed practical notes relevant for replication:
\begin{itemize}
  \item USB port topology affects camera enumeration stability.
  \item Identical camera names in Linux require by-path or by-id mapping.
  \item Mechanical mounting with tape is insufficient for long-term calibration stability.
  \item Uneven wall surfaces can distort tag planarity assumptions.
\end{itemize}

These details strongly influence geometric quality and should be documented as part of method, not as anecdotal notes.



\section{5.23 Calibration Session Documentation Template}
Each calibration session should record:
\begin{itemize}
  \item date/time and operator,
  \item active camera mapping,
  \item board/tag assets used,
  \item command lines executed,
  \item produced calibration files,
  \item overlay validation verdict,
  \item notes on anomalies.
\end{itemize}

This template enables rapid rollback and forensics when a later session behaves unexpectedly.

\section{5.24 Handling Corrupted or Truncated Video Files}
During long recordings, occasional container warnings (e.g., premature end in MKV) were observed. Mitigation strategies:
\begin{itemize}
  \item prefer robust codecs/containers for intermediate work,
  \item validate frame counts immediately after recording,
  \item keep raw and aligned versions separate,
  \item avoid deleting source clips until processing passes sanity checks.
\end{itemize}

\section{5.25 Smoothing Policy for Scientific Reporting}
A key reporting principle in this project: always preserve a raw-processing output and report it explicitly. Smoothed outputs may be used for visualization and operator interpretability, but should not replace raw metrics in scientific claims.

This policy prevents accidental inflation of perceived system quality.

\section{5.26 Camera Placement Engineering}
Camera placement should maximize:
\begin{itemize}
  \item overlap in central operational volume,
  \item diversity of view angles for depth conditioning,
  \item visibility of key calibration tags.
\end{itemize}

The project learned that rotating side cameras too far toward one wall reduces overlap for certain regions and increases triangulation instability.

\section{5.27 Tag Maintenance as a Technical Task}
Marker maintenance was operationally significant. Dirty or damaged tags reduced detection quality and increased reprojection residuals. Routine maintenance should include:
\begin{itemize}
  \item cleaning marker surfaces,
  \item replacing warped paper,
  \item checking tape adhesion and planarity,
  \item updating coordinates after any physical relocation.
\end{itemize}

\section{5.28 Continuous Integration Possibility for Scripts}
Although not yet implemented, scripts can be CI-checked for:
\begin{itemize}
  \item command-line argument integrity,
  \item JSON schema consistency,
  \item report generation consistency on sample clips.
\end{itemize}

This would further improve long-term maintainability.

\section{5.29 Why Some Legacy Files Were Retained Unmodified}
Certain legacy files were retained as historical references instead of being refactored into the new pipeline. This was a deliberate decision to preserve provenance and avoid introducing regressions in already validated scripts.

\section{5.30 Chapter Summary}
Implementation maturity was achieved through repeated practical corrections: better camera mapping, stronger calibration discipline, deterministic session structure, and protocolized reporting. These engineering choices are central to the final performance.
